\chapter{Conclusion}
\label{chapter:conclusion}

The search presented in this thesis is the first ATLAS search dedicated to high-mass $U_1$ leptoquarks in the $\tau\nu(+b)$ final state with a hadronically decaying $\tau$ and missing transverse momentum.
It uses $140~\ifb$ of $pp$ collision data recorded during Run~2 at $\sqrt{s}=13~\TeV$.
The search is motivated by the $\RD$ anomalies, which can be explained by a $U_1$ contribution to $b\to c\tau\nu$.

Both resonant production and non-resonant $t$-channel exchange are probed, extending the sensitivity across different coupling scenarios.
For each topology, two SRs are defined by requiring one $b$ jet or applying a $b$-jet veto, giving four SRs in total to gain better sensitivity.
The optimisation is performed for $\bTagSR$ using a grid search for maximum sensitivity.
The fake-$\hadtau$ background, mainly from $\znunu$+jets, is small. The multi-jet background is negligible after the full SR selections.
The remaining dominant $\wjets$ and top-quark backgrounds are estimated using NFs measured in the corresponding CRs, and their extrapolation is tested in dedicated VRs.

As a result, no significant excess above the SM prediction is observed.
A profile-likelihood fit including statistical and systematic uncertainties is performed to extract the limit on the signal strength $\mu$.
Upper limits are set at 95\% CL on the couplings for $\MU$ between $1.5$ and $3.0~\TeV$ in the purely left-handed scenario.
The limits are obtained using the CL$_{\mathrm{s}}$ method with pseudo-experiments.
At $\MU=1.5~\TeV$, the observed (expected) upper limits on $\gU$ are $1.13$ ($1.11$) for $\beta_L^{23}=1.0$ and $0.52$ ($0.53$) for $\beta_L^{23}=2.2$.
At $\MU=2.5~\TeV$, the corresponding limits are $2.30$ ($2.19$) and $1.22$ ($1.18$).

A supplementary interpretation in the right-handed scenario in $\bTagSR$ is also performed using the asymptotic approximation, excluding the interference effect.
At $\MU=1.5~\TeV$, the observed (expected) upper limits on $\gU$ are $0.84$ ($0.80$) for $\beta_L^{23}=1.0$ and $0.53$ ($0.51$) for $\beta_L^{23}=2.2$.
At $\MU=3.0~\TeV$, the corresponding limits are $1.98$ ($1.96$) and $1.30$ ($1.28$).


These results probe large values of $\beta_L^{23}$ in a region of parameter space not covered by previous $\tau\tau$ searches.
The $\bTagSR$ regions provide most of the sensitivity, while $\bVetoSR$ improves the sensitivity when the $s\tau$ contribution is enhanced at large $\beta_L^{23}$.
The $\tau\nu(+b)$ channel therefore complements the existing searches.

Further improvements are expected from larger datasets and a dedicated optimisation of $\bVetoSR$. Especially the gain from the optimisation is expected to be significant in the low-mass region.